Recent medical EHR agent systems have equipped LLMs with user-designed tools for retrieving semi-structured clinical data, enabling more accurate and reliable decision-making over complex patient timelines. However, curating these tool libraries requires sustained expert effort, long tool chains incur substantial token overhead, and many clinically meaningful questions remain ill-defined in zero-shot settings due to variability across hospitals and guidelines. We introduce \ours{}, an agentic benchmark built on MIMIC-IV comprising 65k+ instances spanning 260 distinct tasks over 63 clinical concepts. Rather than relying on pre-specified functions, models are encouraged to synthesize modular, reusable tools that capture clinical concepts and compose them to answer complex queries. The benchmark is further stratified by \emph{compositional dependency depth} and includes a \emph{longitudinal} surveillance task that tracks patient state across ICU admissions, making reliable zero-shot resolution incredibly challenging. We show that an offline autoformalization pipeline, which translates clinical guidelines into QA-verified function libraries, substantially outperforms zero-shot code generation, improves generalization across patient populations, and reduces per-query token usage by over 20\%.